\section{Conclusion}
\label{sec:conclusion}

This project produced a credible lossless compression tool and a clear comparison between three different design choices. In the final benchmark set, \texttt{v6} achieved the best statistical compression ratio, \texttt{v8} was the fastest version, and \texttt{v5} remained the best overall balance. It combined good compression performance with much lower cost than \texttt{v6} and much better output size than \texttt{v8}.

A main lesson from the project is that compression performance did not come from one isolated algorithmic idea, but from how preprocessing, modeling, and implementation details worked together. BWT and MTF were useful because they made the statistical model more effective, while choices such as block handling, escape estimation, and candidate selection also had a strong impact. The comparison with \texttt{v8} also made the trade-off very clear: higher speed is possible, but usually at the cost of weaker compression.

Perhaps the most interesting result came from file D. Every standard tool stores it at full size because its byte statistics look perfectly random. But \texttt{v10} recognizes that the file was generated by glibc's \texttt{rand()} and compresses it to just 14 bytes by storing the seed instead of the data. This is a practical demonstration of the difference between Shannon entropy and Kolmogorov complexity: the statistics say the file is incompressible, but a short program can reproduce it entirely. For a course on Algorithmic Information Theory, this felt like the right problem to solve.

Future work could extend the PRNG detection to cover other generator families (Mersenne Twister, xorshift, LCG variants), which would broaden the class of high-entropy inputs that \texttt{v10} can handle. On the statistical side, the main gaps are files F and H, where bzip2 and xz still win---better handling of binary-heavy data and longer-range dependencies would help close that gap.
